\chapter{多时间尺度学习：截断信用、延迟更新与异步控制}
\label{app:adiabatic_backprop}

本附录讨论一个可以直接验证的工程问题：当智能系统同时包含毫秒级反应、秒级规划与更慢的模型修订时，怎样在有限显存、算力和通信预算下分配信用，而不把所有状态都锁进一条无限增长的反向计算图。我们把旧稿的“绝热”保留为时间尺度分离的助记词，但不把它当作物理定律或正确性的证明。这里需要区分三件事：截断梯度是一种估计选择，延迟更新是一种调度策略，异步控制是一种运行时架构；三者可以组合，却不彼此推出。HSF-HD在这一层提供一般状态动力学的描述框架，AgentOS中的 $h(t)$、$E/\Theta$、$\Psi$、SPC、TCE、CCM、Boundary、Offloading 与外部记忆 $M_e$ 只是可审计实现的一组接口，不是所有智能系统必须采用的唯一机制。

\section{问题定义：全局反向传播的资源与统计边界}
\label{sec:app_physical_dilemma}

先把对象和假设写清。设离散时间索引为 $k\in\{0,\ldots,T\}$，采样间隔为 $\Delta t>0$ 秒；系统包含快、中、慢三个可微模块，其状态分别为 $x_k^f\in\mathbb{R}^{d_f}$、$x_k^m\in\mathbb{R}^{d_m}$ 与 $x_k^s\in\mathbb{R}^{d_s}$。这些向量默认是经过标度约定的数值表示，不是焦耳能量。输入 $u_k\in\mathcal U$、目标版本 $g_k\in\mathcal G$、结构版本 $A_k\in\mathcal A$ 和有效掩码 $m_k\in\{0,1\}$ 都可能随时间变化。参数集合记为 $\vartheta=(\vartheta_f,\vartheta_m,\vartheta_s)$，单步更新写成
\begin{equation}
  x_{k+1}=F_{\vartheta,k}(x_k,u_k,g_k,A_k,m_k),
  \qquad
  \mathcal J_T(\vartheta)=\sum_{k=1}^{T}w_k\ell_k(x_k,g_k)+\Omega(\vartheta),
\end{equation}
其中 $x_k=(x_k^f,x_k^m,x_k^s)$，$w_k\geq 0$ 为无量纲评估权重，$\ell_k$ 的单位由具体任务决定，$\Omega$ 是正则项。若交叉熵、时延、通信字节和焦耳能耗同时进入目标，必须先分别归一化并报告各分量，不能把它们都称作“能量”后直接相加。目标、结构和掩码随 $k$ 变化时，梯度也依赖这些版本；把它们冻结只是局部分析假设，而不是全程事实。

完整的随时间反向传播对参数的梯度可写为
\begin{equation}
 \nabla_{\vartheta}\mathcal J_T
 =\sum_{k=1}^{T}\left[
 \frac{\partial \ell_k}{\partial \vartheta}
 +\frac{\partial \ell_k}{\partial x_k}
 \sum_{j=0}^{k-1}
 \left(\prod_{r=j+1}^{k-1}\frac{\partial F_{\vartheta,r}}{\partial x_r}\right)
 \frac{\partial F_{\vartheta,j}}{\partial \vartheta}
 \right]+\nabla_{\vartheta}\Omega .
\end{equation}
这是链式法则的离散展开；它不假设信息以无限速度传播，也不意味着模型与物理系统整体同构。真正的工程困难来自保存或重算中间激活、长雅可比乘积的数值条件，以及非平稳数据使早期梯度对当前参数失去代表性。若每步保存的激活占用 $a_k$ 字节，朴素实现的激活显存约为 $M_{act}=\sum_{k=1}^{T}a_k$；只有在 $a_k\approx a$ 时才可简写为 $\mathcal O(Ta)$。检查点重算可以用额外计算换取较低显存，隐式梯度、可逆层或前向模式又有不同边界，所以“完整BPTT必然爆显存”不是无条件数学结论。

长乘积还带来梯度消失或爆炸。若在一个固定版本窗口内有 $\left\lVert \partial F_{\vartheta,r}/\partial x_r\right\rVert_2\leq\rho$，则从第 $j$ 步到第 $k$ 步的敏感度上界为 $\rho^{k-j-1}$。当 $\rho<1$ 时远期影响衰减，当 $\rho>1$ 时上界增长；但该界只给出范数风险，不能说明某个方向上的真实梯度，也不能把数值爆炸解释为热力学熔断。归一化、门控、残差连接、谱约束和梯度裁剪可以改变数值条件，却不自动保证任务成功。反过来，稳定的梯度也可能对应错误目标、错误数据或权限越界，因此数值稳定、统计有效与安全合规需要分别验收。

多时间尺度系统还面临信用混叠。传感噪声、短期控制误差、任务切换和真正的慢层模型偏差都可能在同一损失中出现；如果每一步都让慢参数响应，慢层会追逐高频波动。可这不是因为噪声必然“倒灌”，而是因为估计器的频率响应、采样相关性和更新规则没有被设计。设快层每步更新，慢层每 $N_s$ 步更新一次。简单求和的慢层梯度估计 $G_q^s=\sum_{k=qN_s}^{(q+1)N_s-1}g_k^s$ 会随窗口长度改变尺度；使用均值 $\bar G_q^s=N_s^{-1}G_q^s$、指数加权或稳健截尾会产生不同偏差与方差。若梯度具有持续同号分量，累积会增强它；若梯度正负交替，可能抵消；若噪声有偏或长程相关，抵消并不成立。因而不能仅凭“多步相加”宣称高频噪声已经消失，必须用功率谱、窗口消融、估计方差和任务指标验证。

本章采用的工程目标是受预算的信用分配，而不是模拟某种宇宙过程。对任一方案，我们至少记录峰值显存、每步计算量、墙钟时延、梯度范数与方差、长期任务回报、切换后的恢复时间及罕见事件漏检率；对异步实现还要记录状态版本陈旧度、丢包、乱序和回退次数。基线包括完整BPTT、固定窗口截断、只训练快层、周期更新慢层和检查点重算。若时间尺度分离只减少显存却显著损害长期信用，它只是一项资源折中；若训练损失下降而部署回报没有改善，也不能称为学习成功。只有在相同数据、参数量、硬件和停止规则下同时报告质量与资源前沿，后续的截断、累积和异步接口才具有可比较的意义。
